\section{Methods and Benchmark Design}
\label{sec:methods}

% Question: What do the methods of this project consist of?
% Answer: Two connected parts: NOVA, the detector-level inference method,
% and the injector, which builds the truth-known benchmark on which NOVA and
% extraction-based analyses are compared.
% Evidence: The structure of this section.

The methods comprise two parts. The first is NOVA, the detector-level
inference method: a forward model of the calibrated detector time series
whose parameters are fitted directly to the retained pixels. The second is
the injector, which constructs the benchmark: a detector cube built around
real out-of-transit JWST data in which a known transmission spectrum has
been hidden, so that recovery accuracy can be measured absolutely. The two
parts are developed and operated separately. Recovery never sees the
injected spectrum, generator-only products, or scoring code until its
prediction is closed, and no calibration, mask, response, threshold,
starting point, or convergence decision is allowed to use the injected
truth or any recovered-depth residual. The specific NOVA implementation
evaluated in this assessment is its frozen science release, NOVA-S: an
immutable, hash-authenticated snapshot of code, configuration, and
calibration products. Later refinements are candidates or diagnostics and
are not part of NOVA-S. Model classes that failed the truth-blind
development gates are recorded rather than discarded
(Appendix~\ref{app:nova}).

\subsection{\nova{}}

% Question: What problem does NOVA solve, and in what form?
% Answer: It inverts a forward model of the detector: the same physical
% quantities that generate the pixels are inferred from them in one
% likelihood.
% Evidence: Sections 1.3 and 1.4.

NOVA treats spectral recovery as one inverse problem. A forward model maps
the physical description of the observation, the transmission spectrum
$D(\lambda)$, the limb darkening $u(\lambda)$, the stellar continuum, the
spatial response, and the additive background, to a prediction for every
retained detector sample. The inverse problem infers those quantities from
the observed samples. Nothing is extracted, binned, or fitted in a
separate downstream stage: the wavelength-dependent transit is estimated
in the same likelihood that describes the detector.

\subsubsection{The detector forward model}

% Question: What does NOVA predict for each detector sample?
% Answer: Continuum times exposure-integrated transit times spatial
% response, corrected by a fixed source-curvature factor, plus an additive
% background.
% Evidence: The handoff's exact method note in the presentation's notation.

The retained detector samples are organised into wavelength groups: one
group $g$ collects the pixels $P_g$ of one spectral order that share one
native detector column, and $g=g(p)$ denotes the group of pixel $p$. For
integration $t$ and pixel $p$, NOVA predicts
\begin{equation}
 \widehat Y_{tp}
 = \gamma_t\, C_{tg}\,
   \mathcal{T}_{tg}(D_g, \bm{u}_g;\, \Omega)\,
   q^\star_{pg}
 + B_{tp},
 \label{eq:nova-forward}
\end{equation}
and, where order supports overlap, the source term is summed over the
groups $\mathcal{G}(p)$ that contribute to the pixel before the residual
is formed. Reading outward from the star: the continuum $C_{tg}$ sets the
brightness of the group at time $t$; the exposure-integrated transit
factor $\mathcal{T}_{tg}$ dims that light according to the depth $D_g$ and
limb darkening $\bm{u}_g$ at the group's wavelength; the fixed factor
$\gamma_t$ applies a small correction for slow source variation over the
visit; the spatial response $q^\star_{pg}$ distributes the group's light
across its pixels, with $\sum_{p\in P_g} q^\star_{pg}=1$; and $B_{tp}$ is
the additive background. The group geometry, wavelength assignment, and
response support are defined by the current public trace and wavelength
calibration \citep{BainesEtAl2023Trace,BainesEtAl2023Wavelength} and the
associated spectral-response kernel
\citep{DarveauBernierEtAl2022,JWSTPipeline2025}. The background is
additive and is never multiplied by the transit or by $\gamma_t$: a change
in transit depth cannot be imitated by dimming the background along with
the star.

The unit of the likelihood is one retained pixel in one integration.
Valid samples have finite values, positive uncertainties, and no
\texttt{DO\_NOT\_USE} flag; 97{,}557 pixels per integration are retained,
organised into 2{,}720 groups after exclusion of six unreliable edge
groups.

\subsubsection{The shared spectrum and transit model}

% Question: How are the depths and the transit generated?
% Answer: One physical spectrum shared by both orders, an exact Keplerian
% transit averaged over each exposure, and small fitted limb-darkening
% deviations around a theory reference.
% Evidence: The exact method note and presentation slides 12 and 13.

Both spectral orders evaluate the same physical spectrum on their own
wavelength grids, $D_g=D(\lambda_g)$. The spectrum is carried by $K=160$
adaptive basis coefficients,
\begin{equation}
  D_j = D_{\mathrm{global}}\,
  \frac{\exp\!\left[(\bm{H}_v\bm{v})_j\right]}
       {\sum_k w_k \exp\!\left[(\bm{H}_v\bm{v})_k\right]} ,
  \label{eq:depth-basis}
\end{equation}
which separates the achromatic depth $D_{\mathrm{global}}$ from a
weighted-zero-mean chromatic morphology and keeps all depths positive. A
two-coordinate order-discrepancy term $\Delta$ exists in the model, with
order 1 seeing $D(\lambda)+\Delta$ and order 2 seeing $D(\lambda)-\Delta$;
in NOVA-S it is pinned to numerical zero and carries no physical freedom.
No atmospheric template, wavelength-local smoother, or truth-informed
spectral prior is used.

The orbital state
$\Omega=(P,\,a/R_\star,\,e,\,\omega_{\mathrm{peri}},\,b,\,t_0)$ is fixed
at the values listed in Appendix~\ref{app:nova}. The instantaneous
limb-darkened transit factor is computed with an exact Keplerian engine
\citep{Jaxoplanet2025}, and $\mathcal{T}_{tg}$ is its average over the
finite exposure, evaluated with 21 weighted model evaluations per
integration. Quadratic limb darkening \citep{MandelAgol2002} uses the
bounded coordinates of \citet{Kipping2013},
\begin{equation}
 u_1=2\sqrt{q_1}\,q_2, \qquad
 u_2=\sqrt{q_1}\,(1-2q_2),
 \label{eq:kipping-coordinates}
\end{equation}
with the wavelength-dependent theory reference taken from the STAGGER
grid \citep{MagicEtAl2015}. The fitted freedom around that reference is
deliberately small: one offset and one normalised log-wavelength slope per
transformed coordinate and order, eight nonlinear coordinates in total,
constrained by theory-deviation and curvature penalties.

\subsubsection{The empirical spatial response}

% Question: Where does q-star come from?
% Answer: The transit itself reveals where the starlight lands: a robust
% per-pixel regression of the data on a depth-free transit template, with
% amplitudes normalised within each group.
% Evidence: Presentation slide 14 and the exact method note.

The spatial response is estimated from the delivered cube itself, using
the transit as its own probe: starlight dims during transit, the additive
background does not, so the transit-correlated amplitude of each pixel
measures how much starlight lands there. For each order a depth-free
event template $\phi_o(t)$ is built from a robust median of
baseline-normalised bright-pixel residuals, and every qualifying pixel is
fitted robustly as
\begin{equation}
 Y_{tp} = A_{pg}\,\phi_{o}(t) + m_p + n_p\,\tau_t + \varepsilon_{tp},
 \label{eq:qstar-fit}
\end{equation}
with a level $m_p$, a linear drift $n_p\tau_t$, and the transit-correlated
amplitude $A_{pg}$. The response is the normalised amplitude,
$q^\star_{pg}=A_{pg}/\sum_{p'\in P_g}A_{p'g}$, so total group brightness
remains with the fitted continuum. Alternating folds, stratified across
the transit phases, provide independent estimates, and the response is
admitted only after cross-fold flux-distribution and centroid stability
gates (Appendix~\ref{app:nova}). The frozen response carries one
documented limitation: its smoother operates on an integer trace-relative
coordinate, which produces centroid resets and a small amplitude mismatch
at the reddest columns. Because the continuum fixes each group's
brightness while $q^\star$ fixes its spatial profile, such a response
over-amplitude acts like a dilution and makes the recovered transit
shallower at the affected columns. The defect is documented rather than
silently repaired; its causal assessment belongs to the Results.

\subsubsection{The continuum and the additive background}

% Question: How are the stellar continuum and the background modelled and
% calibrated?
% Answer: A level and slope per group for the continuum; eight spatial maps
% crossed with two temporal patterns for the background, selected and
% anchored on off-trace pixels only.
% Evidence: Presentation slide 15 and the exact method note.

Each group's continuum is a level and a linear time slope,
$C_{tg}=\beta_{0g}+\beta_{1g}\tau_t$, giving 5{,}440 linear coefficients.
The additive background is a sum of eight fixed spatial maps $\Phi_{pk}$
with two temporal patterns,
\begin{equation}
 B_{tp}=\sum_{k=1}^{8}\Phi_{pk}
 \left[\alpha_{k0}+\alpha_{k1}\,G_1(t)\right],
 \label{eq:background-model}
\end{equation}
sixteen coefficients $\alpha_{km}$ in all. Everything about this model is
calibrated on pixels away from the stellar traces. The spatial family and
rank were selected by hold-out prediction on off-trace pixels; the
per-integration amplitudes of the selected maps were decomposed into
principal temporal components, and $G_1(t)$ is the first component whose
absolute correlation with a fixed transit-reference template does not
exceed 0.25, so the background cannot quietly learn the transit. An
off-trace weighted least-squares estimate provides the expectation
$\bm{\alpha}_{\mathrm{off}}$ and its precision $\bm{N}_{\mathrm{off}}$,
which enter the objective as a data-derived anchor: the background is
anchored, not fixed, and is fitted jointly with the continuum at every
step. Recovery never subtracts a fixed background. The likelihood contains
no Gaussian-process term for time-correlated residuals; the robust weights
below handle outliers, and residual temporal correlation is tracked as a
diagnostic.

\subsubsection{The source-curvature factor and detector uncertainties}

% Question: What are gamma_t and sigma_tp?
% Answer: A fixed out-of-transit curvature correction selected on held-out
% folds, and per-sample uncertainties scaled per order and per detector row
% from out-of-transit cross-validation.
% Evidence: The exact method note, Sections 8 and 10.

The factor $\gamma_t$ corrects for slow variation of the source over the
visit. It was calibrated on the out-of-transit integrations only, split
into two folds: stable trace flux was integrated per order, the ratio of a
weighted quadratic to its affine approximation defined each candidate, and
null, common, and per-order candidates were ranked by held-out
standardised chi-squared. The common candidate was selected because the
per-order candidate failed the requirement that both orders improve
individually; the selection scores are in Appendix~\ref{app:nova}.
$\gamma_t$ has unit out-of-transit mean, multiplies only the stellar and
transit term of Equation~\ref{eq:nova-forward}, and adds no fitted
parameters.

The per-sample uncertainties are
$\sigma_{tp}=\mathrm{ERR}_{tp}\,N_o\,s_{r(p)}$: the calibrated ERR array,
scaled once per order by $N_o$, and once per detector row by a guard
$s_{r(p)}\geq 1$ derived from bidirectional out-of-transit
cross-validation, so rows whose quiet-time residuals exceed their nominal
uncertainties are down-weighted. No wavelength information, injected
truth, or recovered spectrum enters either scale.

\subsubsection{The complete objective}

% Question: What does NOVA minimise?
% Answer: A robust detector mismatch over the retained mask plus the
% off-trace background anchor and the limb-darkening and order-discrepancy
% penalties.
% Evidence: Presentation slide 16.

With the standardised residual
\begin{equation}
 r_{tp}=\frac{Y_{tp}-\widehat Y_{tp}
 (D,\bm{u},\bm{\beta},\bm{\alpha},\Delta;\,q^\star,\Omega)}{\sigma_{tp}},
 \label{eq:residual}
\end{equation}
where the semicolon separates fitted parameters from frozen inputs, the
estimate is
\begin{equation}
\begin{split}
 &(\widehat D,\widehat{\bm{u}},\widehat{\bm{\beta}},
  \widehat{\bm{\alpha}},\widehat\Delta)
 = \arg\min\Big\{ \\
 &\quad\sum_{(t,p)\in\mathcal{M}}\rho_{1.345}(r_{tp}) \\
 &\quad+\tfrac{1}{2}(\bm{\alpha}-\bm{\alpha}_{\mathrm{off}})^{\mathsf T}
   \bm{N}_{\mathrm{off}}(\bm{\alpha}-\bm{\alpha}_{\mathrm{off}}) \\
 &\quad+\tfrac{1}{2}\|\bm{r}_{\mathrm{LD,theory}}(\bm{u})\|_2^2
  +\tfrac{1}{2}\|\bm{r}_{\mathrm{LD,curv}}(\bm{u})\|_2^2 \\
 &\quad+\tfrac{1}{2}\|\bm{r}_{\mathrm{order}}(\Delta)\|_2^2
 \Big\},
\end{split}
 \label{eq:objective}
\end{equation}
where $\mathcal{M}$ is the retained sample mask and $\rho_{1.345}$ is the
Huber loss. The four terms are the detector mismatch, the off-trace
background anchor, the limb-darkening theory and smoothness penalties, and
the order-discrepancy penalty that pins $\Delta$ in NOVA-S.

\subsubsection{The nested solver}

% Question: How is the objective minimised, and what counts as convergence?
% Answer: Robust weights outside, bounded trust-region steps in the middle,
% and an exact linear solve for continuum and background inside; acceptance
% requires the preregistered gates and agreement of two independent starts.
% Evidence: Presentation slides 17 and 18, and the frozen solver code.

The parameters split into a nonlinear block
$\bm{\theta}=(D,\bm{u},\Delta)$ and a linear block
$(\bm{\beta},\bm{\alpha})$, and the solver nests three layers. Innermost,
at fixed $\bm{\theta}$ and fixed robust weights, variable projection
\citep{GolubPereyra1973} solves the linear problem exactly,
\begin{equation}
\begin{split}
 (\widehat{\bm{\beta}},\widehat{\bm{\alpha}})=
 \arg\min_{\bm{\beta},\bm{\alpha}}\;
 &\frac{1}{2}\sum_{t,p}\frac{w_{tp}}{\sigma_{tp}^2}
 \left[Y_{tp}-\widehat Y_{tp}(\bm{\theta},\bm{\beta},\bm{\alpha})\right]^2 \\
 &+\frac{1}{2}(\bm{\alpha}-\bm{\alpha}_{\mathrm{off}})^{\mathsf T}
   \bm{N}_{\mathrm{off}}(\bm{\alpha}-\bm{\alpha}_{\mathrm{off}}),
\end{split}
 \label{eq:varpro}
\end{equation}
by factorising the sparse continuum block and reducing the background
coupling to a sixteen-dimensional Schur complement; the profiled residual
and Jacobian differentiate through this exact inner optimum
(Appendix~\ref{app:nova}). In the middle, bounded trust-region reflective
least squares \citep{BranchColemanLi1999} minimises the profiled
objective over $\bm{\theta}$ in scaled coordinates, with stopping
controls $\mathrm{gtol}=10^{-6}$, $\mathrm{ftol}=10^{-9}$,
$\mathrm{xtol}=10^{-8}$, and at most 150 function evaluations per solve.
Outermost, exact Huber iteratively reweighted least squares
\citep{Huber1964} updates the weights,
$w_{tp}=\min(1,\,1.345/|e_{tp}|)$ for standardised residual $e_{tp}$,
holding them fixed during each trust-region solve, for at most 50 cycles.

Acceptance is stricter than the optimiser's own termination. A fit is
accepted only when the robust weights (maximum change below $0.001$), the
objective (relative change below $10^{-8}$), and the shared spectrum
(root-mean-square changes below 0.01 and 0.05 ppm on the two spectral
summaries) have stopped moving, the scale-aware Karush-Kuhn-Tucker and
first-order measures are below $10^{-6}$ and $0.002$, a deterministic
material-descent audit finds no remaining feasible improvement
(Appendix~\ref{app:nova}), and a final confirmation fit at the recomputed
terminal weights succeeds. Two deterministic starts are run independently
with no state transfer, one at the recorded white-light mean depth with
zero morphology and theory limb darkening, the other with a fixed
zero-mean 250 ppm morphology perturbation and small fixed limb-darkening
offsets; they must agree within 0.1 and 0.5 ppm on the two spectral
summaries and within $10^{-8}$ in relative objective. Among physical
starts the executable rule prefers a converged start over a merely valid
one and then selects the lowest truth-blind robust objective. Both starts
of the frozen release converged, at robust cycle 14; truth error plays no
role in any of these decisions.

\subsection{The injector}

% Question: How is a truth-known benchmark built from real JWST data?
% Answer: A known transit deficit, applied only to a denoised target-source
% expectation, is added to real out-of-transit integrations, so background
% and noise are real and the truth is known.
% Evidence: Presentation slide 20 and the injector record.

The injector answers the validation problem of
Section~\ref{sec:detector-to-validation}: real observations provide no
ground truth, so a truth is hidden inside real data. The baseline is the
observed out-of-transit detector time series of the WASP-17 NIRISS/SOSS
visit, 229 calibrated integrations that already contain the target star,
the additive background, field sources, and the realised noise. A known
transmission spectrum $D^{\mathrm{inj}}(\lambda)$ defines per-order
transit factors $\bm{T}^{\mathrm{true}}_{o,t}$, and the delivered cube is
\begin{equation}
\begin{split}
 Y^{\mathrm{inj}}_{tp}
 = {}& Y^{\mathrm{OOT}}_{tp} \\
 & + \sum_{o}\left[\bm{A}_o\!\left\{\bm{s}_{o,t}\odot
   \left(\bm{T}^{\mathrm{true}}_{o,t}-\bm{1}\right)\right\}\right]_p ,
\end{split}
 \label{eq:v45-injector}
\end{equation}
where $\bm{s}_{o,t}$ is a denoised expectation of the target source for
order $o$ and $\bm{A}_o$ is the corresponding detector projection built
from the same public instrument references
\citep{BainesEtAl2023Trace,BainesEtAl2023Wavelength,DarveauBernierEtAl2022}.
Only the target-source expectation transits: the observed background,
contaminants, and noise are left untouched, and the source expectation is
denoised so no particular noise realisation is dimmed. The support logic
fails closed, and the synthetic event occupies integrations 51--177,
leaving 102 out-of-transit integrations for the truth-blind calibrations
of the previous subsection. The full construction, delivery manifest, and
separation of generator, delivery, recovery, and scoring are specified in
Appendix~\ref{app:benchmarking}.

\subsubsection{Comparison analyses and fairness status}

% Question: Against what is NOVA compared, and on what terms?
% Answer: An authentic Ahsoka analysis of the same cube now, further
% pipelines later; fairness is defined at delivery, and the group-stage
% preprocessing experiment remains open.
% Evidence: The injector record and Louie et al. (2025).

Every recovery method receives the same opaque cube, uncertainties, flags,
time grid, and declared wavelength support; unsupported bins are excluded
symmetrically, and each method keeps its own nuisance model and inference
strategy. The first external comparison is an authentic Ahsoka reduction
following the analysis topology of the published WASP-17~b spectrum
\citep{LouieEtAl2025}. One fairness limitation is open: the principal
comparison pipelines apply their most distinctive $1/f$ correction on
detector groups before ramp fitting, and an integration-level cube
contains no groups. A preregistered paired experiment therefore injects
the same deficit at group stage and at integration stage and runs NOVA and
the comparison pipelines on both; it has not been completed and is not
part of the frozen result. Later rounds will add Eureka!,
supreme-SPOON/exoTEDRF, transitspectroscopy, and the official JWST
pipeline with ATOCA extraction
\citep{BellEtAl2022,Radica2024,Espinoza2022TransitSpectroscopy,%
JWSTPipeline2025,DarveauBernierEtAl2022,LouieEtAl2025}, each with its
authentic truth-blind settings on the same delivered cube.

\subsubsection{Reporting support and evaluation metrics}

% Question: On what wavelength support are methods compared, and with which
% metrics?
% Answer: Whole bins of a common R=100 grid inside the directly calibrated
% wavelength domain, scored by absolute RMSE, mean bias, and morphology
% RMSE, with per-recovery diagnostics now and ensemble coverage tests later.
% Evidence: The exact method note and the benchmark appendix.

Recovered and injected spectra are compared on a common $R=100$ grid,
restricted prospectively to bins whose detector support lies entirely
inside the directly calibrated wavelength domain of the trace solution;
beyond its last training anchor the wavelength solution is a polynomial
extrapolation, so those bins are excluded in full. For this visit the
rule keeps 141 bins, the last ending near $2.740\,\mu\mathrm{m}$
(Appendix~\ref{app:benchmarking}). The rule is fixed across methods,
changes no fit, and is applied only after recovery.

With $e_i=\widehat D_i-D_i^{\mathrm{inj}}$ on the common grid, the scores
are
\begin{equation}
 \mathrm{RMSE}_{\mathrm{abs}}=
 \left(\frac{1}{N}\sum_i e_i^2\right)^{1/2},
 \label{eq:absolute-rmse}
\end{equation}
\begin{equation}
 \mathrm{RMSE}_{\mathrm{morph}}=
 \left[\frac{1}{N}\sum_i(e_i-\bar e)^2\right]^{1/2},
 \label{eq:morphology-rmse}
\end{equation}
so a constant offset is distinguished from an incorrect spectral shape,
together with mean and median bias, feature-amplitude error, and
order-specific residuals. Standardised residuals and the spectral
covariance from the profiled normal system accompany each individual
recovery; empirical coverage of injected depths requires repeated noise
realisations and belongs to the planned ensemble. Detector residuals,
robust weights, rank, conditioning, boundary activity, and convergence
diagnostics accompany every score: a low error is not sufficient if a fit
fails its validity gates.
